\section{Xác suất}

\subsection{Đề 1}
\Opensolutionfile{ans}[ans/ans-0-B1]
\caulc
\begin{ex}%[0D0N2-5]
	Gieo một con súc sắc. Xác suất để mặt chấm chẵn xuất hiện là
	\choice
	{$0{,}2$}
	{$0{,}3$}
	{$0{,}4$}
	{\True $0{,}5$}
	\loigiai{
		Gọi A là biến cố: \lq\lq Mặt chấm chẵn xuất hiện\rq\rq.\\
		Ta có $n(A)=3$, $n(\Omega)=6$\\
		$\Rightarrow \mathrm{P}(A)=0,5$.}
\end{ex}
\begin{ex}%[0D0N2-5]
	Một hội nghị có $15$ nam và $6$ nữ. Chọn ngẫu nhiên $3$ người vào ban tổ chức. Xác suất để $3$ người lấy ra là nam:
	\choice
	{$\dfrac{1}{2}$}
	{\True $\dfrac{91}{266}$}
	{$\dfrac{4}{33}$}
	{$\dfrac{1}{11}$}
	\loigiai{
		$n(\Omega)=\mathrm{C}_{21}^3=1330$\\
		Gọi $\mathrm{A}$ là biến cố: \lq\lq 3 người lấy ra là nam\rq\rq.\\ Khi đó, $n(A)=\mathrm{C}_{15}^3=455$.
		Vậy xác suất để 3 người lấy ra là nam là: $\mathrm{P}(A)=\dfrac{n(A)}{n(\Omega)}=\dfrac{13}{38}=\dfrac{91}{266}$.
	}
\end{ex}
\begin{ex}%[0D0N2-5]
	Từ một hộp chứa 11 quả cầu màu đỏ và 4 quả cầu màu xanh, lấy ngẫu nhiên đồng thời 3 quả cầu. Xác suất để lấy được 3 quả cầu màu xanh bằng
	\choice
	{$\dfrac{4}{165}$}
	{\True $\dfrac{4}{455}$}
	{$\dfrac{33}{91}$}
	{$\dfrac{24}{455}$}
	\loigiai{
		Ta có: $n(\Omega)=\mathrm{C}_{15}^3=455$.
		Gọi $A$ là biến cố \lq\lq lấy được 3 quả cầu màu xanh\rq\rq.\\ Khi đó $n(A)=\mathrm{C}_4^3=4$.\\
		Vây $\mathrm{P}(A)=\dfrac{4}{455}$.
	}
\end{ex}
\begin{ex}%[0D0N2-5]
	Từ các chữ số $1, 2, 4, 6, 8, 9$ lấy ngẫu nhiên mộtt số. Xác suất để lấy được số lẻ bằng
	\choice
	{\True $\dfrac{1}{3}$}
	{$\dfrac{1}{2}$}
	{$\dfrac{1}{4}$}
	{$\dfrac{1}{6}$}
	\loigiai{
		Số phần tử của không gian mẫu là $n(\Omega)=\mathrm{C}_6^1=6$.\\
		Gọi $A$ là biến cố \lq\lq lấy được số lẻ\rq\rq. Số phần tử của biến cố là $n(A)=\mathrm{C}_2^1=2$.\\
		Xác suất của biến cố $A$ là $\mathrm{P}(A)=\dfrac{n(A)}{n(\Omega)}=\dfrac{1}{3}$.}
\end{ex}
\begin{ex}%[0D0N2-5]
	Gieo một đồng xu cân đối đồng chất liên tiếp hai lần. Tính xác suất để cả hai lần gieo đều được mặt sấp.
	\choice
	{\True $\dfrac{1}{4}$}
	{$\dfrac{1}{6}$}
	{$\dfrac{1}{8}$}
	{$\dfrac{1}{2}$}
	\loigiai{
		Gọi $\Omega$ là không gian mẫu. Gieo một đồng xu hai lần liên tiếp nên $n(\Omega)=2\cdot2=4$
		Gọi $A$ \lq\lq  Cả hai lần gieo đều mặt sấp\rq\rq\,nên $n(A)=1.1=1$
		Vậy $\mathrm{P}(A)=\frac{n(A)}{n(\Omega)}=\frac{1}{4}$.}
\end{ex}
\begin{ex}%[0D0N2-5]
	Gieo một con súc sắc cân đối và đồng chất. Xác suất xuất hiện mặt có số chấm chia hết cho 3 bằng
	\choice
	{$\dfrac{1}{3}$}
	{\True $\dfrac{1}{3}$}
	{$\dfrac{2}{3}$}
	{$\dfrac{2}{3}$}
	\loigiai{
		Ta có $n(\Omega)=6$ và $n(A)=2$. \\Vậy $\mathrm{P}(A)=\dfrac{1}{3}$.}
\end{ex}
\begin{ex}%[0D0N2-5]
	Một đoàn đại biểu gồm $5$ người được chọn ra từ một tổ gồm $8$ nam và $7$ nữ để tham dự hội nghị. Xác suất để chọn được đoàn đại biểu có đúng $2$ người nữ là
	\choice
	{\True $\dfrac{56}{143}$}
	{$\dfrac{140}{429}$}
	{$\dfrac{1}{143}$}
	{$\dfrac{28}{715}$}
	\loigiai{
		Số phần tử của không gian mẫu: $n(\Omega)=\mathrm{C}_{15}^5$.\\
		Gọi biến cố $A$ : \lq\lq Chọn được đoàn đại biểu có đúng 2 người nữ\rq\rq\,$\Rightarrow n(A)=\mathrm{C}_7^2 \cdot \mathrm{C}_8^3 $.\\
		Vậy xác suất cần tìm là: $\mathrm{P}(A)=\dfrac{n(A)}{n(\Omega)}=\dfrac{56}{143}$.}
\end{ex}
\begin{ex}%[0D0N2-5]
	Trong trò chơi \lq\lq Hãy chọn giá đúng\rq\rq\,chiếc kim của bánh xe có thể dừng lại ở 1 trong 20 nấc điểm với khả năng như nhau. Tính xác xuất để trong hai lần quay, chiếc kim của bánh xe đó dừng lại ở hai nấc điểm khác nhau.
	\choice
	{$\dfrac{1}{20}$}
	{\True $\dfrac{19}{20}$}
	{$\dfrac{1}{10}$}
	{$\dfrac{9}{10}$}
	\loigiai{
		Số phần tử của không gian mẫu: $n(\Omega)=20^2=400$.\\
		Gọi $A$ là biến cố: \lq\lq trong hai lần quay, chiếc kim của bánh xe đó dừng lại ở hai nấc điểm khác nhau\rq\rq. \\
		Ta có $n(A)=\mathrm{A}_{20}^2=380$.\\Vậy $\mathrm{P}(A)=\dfrac{n(A)}{n(\Omega)}=\dfrac{380}{400}=\dfrac{19}{20}$.}
\end{ex}
\begin{ex}%[0D0N2-5]
	Lấy ngẫu nhiên 3 bi. Tính xác suất để 3 bi lấy ra có ít nhất một bi đỏ.
	\choice
	{$\dfrac{3}{4}$}
	{$\dfrac{10}{21}$}
	{$\dfrac{2}{7}$}
	{\True $\dfrac{37}{42}$}				
	\loigiai{				
		Chọn ngẫu nhiên 3 bi trong 9 bi có $n(\Omega)=\mathrm{C}_9^3=84$.\\
		Chọn 3 bi trong đó có ít nhất 1 bi đỏ là: $n(A)=\mathrm{C}_4^3+\mathrm{C}_4^2 \mathrm{C}_5^1+\mathrm{C}_4^1 \mathrm{C}_5^2=74$. Xác suất để 3 bi được chọn có ít nhất 1 bi đỏ là: $\mathrm{P}(A)=\dfrac{n(A)}{n(\Omega)}=\dfrac{74}{84}=\dfrac{37}{42}$.}
\end{ex}				
\begin{ex}%[0D0N2-5]
	Xác suất để tổng số chấm xuất hiện trên hai mặt của hai con súc sắc bằng 7 là
	\choice
	{$\dfrac{6}{7}$}
	{$\dfrac{1}{7}$}
	{$\dfrac{1}{6}$}
	{\True $\dfrac{5}{6}$}
	\loigiai{
		Số phần tử của không gian mẫu là $n(\Omega)=6 \cdot 6=36$.\\
		Gọi $A$ là biến cố \lq\lq Tổng số chấm xuất hiện trên hai mặt của hai con súc sắc bằng 7\rq\rq.\\ 
		Ta có$A=\lbrace(1 ; 6),(2 ; 5),(3 ; 4),(4 ; 3),(5 ; 2),(6 ; 1)\rbrace \Rightarrow n(A)=6$.\\
		Vậy $\mathrm{P}(A)=\dfrac{6}{36}=\dfrac{1}{6}$.}
\end{ex}
\begin{ex}%[0D0N2-5]
	Một tổ có $6$ nam và $4$ nữ. Chọn ngẫu nhiên hai người. Tính xác suất sao cho trong hai người được chọn có ít nhất một người là nữ.
	\choice
	{$\dfrac{4}{5}$}
	{\True $\dfrac{2}{3}$}
	{$\dfrac{2}{15}$}
	{$\dfrac{1}{3}$}
	\loigiai{	
		Số phần tử không gian mẫu là số cách chọn 2 người trong 10 người $n(\Omega)= \mathrm{C}_{10}^2$.\\
		Gọi $A$ là biến cố hai người có ít nhất một người là nữ.\\
		Suy ra $\overline{A}$ là biến cố \lq\lq cả hai người là nam\rq\rq. Ta có $n(\overline{A})=\mathrm{C}_6^2$. Do đó xác suất $\mathrm{P}(A)=1-\dfrac{\mathrm{C}_6^2}{\mathrm{C}_{10}^2}=\dfrac{2}{3}$.}
\end{ex}
\begin{ex}%[0D0N2-5]
	Một hộp đựng $5$ viên bi đỏ, $4$ viên bi xanh. Lấy ngẫu nhiên $3$ viên bi từ hộp đó. Tính xác suất lấy được ít nhất một viên đỏ.
	\choice
	{$\dfrac{1}{21}$}
	{$\dfrac{37}{42}$}
	{$\dfrac{5}{42}$}
	{\True $\dfrac{20}{21}$}
	\loigiai{				
		Ta có $n(\Omega)=\mathrm{C}_9^3=84$.\\
		Gọi $A$ là biến cố cần tìm. \\
		Ta có $n(A)= \mathrm{C}_5^3+\mathrm{C}_5^2 \cdot \mathrm{C}_4^1+\mathrm{C}_5^1 \cdot \mathrm{C}_4^2=80$.\\
		Khi đó $\mathrm{P}(A)=\dfrac{80}{84}=\dfrac{20}{21}$.}
\end{ex}						
\Closesolutionfile{ans}
\indapan{6}{ans/ans-0-B1}
\cauds %4cau
\Opensolutionfile{ans}[ans/ans-0-B1-DS]
\begin{ex} %[0D0V2-5]
	Trong lớp $10 \mathrm{~A}$ có 25 bạn nam và 21 bạn nữ. Giáo viên chon ngẫu nhiên 3 bạn trong lớp để làm cán bô lớp. Khi đó:
	\choiceTF
	{\True Số cách chon ra 3 bạn trong lớp 10A là 15180 (cách)}
	{\True Xác suất của các biến cố \lq\lq Ba bạn được chon đều là nam\rq\rq\,bằng: $\dfrac{5}{33}$}
	{Xác suất của các biến cố \lq\lq Ba bạn được chon đều là nữ\rq\rq\,bằng: $\dfrac{133}{1158}$}
	{\True Xác suất của các biến cố \lq\lq Trong ba hoc sinh được chon có hai bạn nam và môt bạn nữ\rq\rq\,bằng: $\dfrac{105}{253}$}
	\loigiai{
		\begin{itemchoice}
			\itemch Đúng. Số cách chọn ra 3 bạn trong lớp 10A gồm 46 bạn ( 25 bạn nam và 21 bạn nữ) là: $\mathrm{C}_{46}^3=15180$ (cách). Do đó, $n(\Omega)=15180$.
			\itemch Đúng. Suy ra $n(A) = 2300$.
			\itemch Sai. Số cách chọn được 3 bạn nữ từ 21 bạn nữ là: $\mathrm{C}_{21}^3=1330$ (cách).
			\itemch Đúng. Số cách chọn được 2 bạn nam và 1 bạn nữ là: $\mathrm{C}_{25}^2 \cdot \mathrm{C}_{21}^1=6300$ (cách).
		\end{itemchoice}
	}
\end{ex}
%ex_2
\begin{ex}%[0D0V2-5]
	Lớp $10 B$ có 40 hoc sinh, trong đó có nhóm siêu quậy gồm Viêt, Đức, Cường, Thịnh. Cô giáo goi ngẫu nhiên 2 bạn trong lớp để kiểm tra bài cũ. Khi đó:
	\choiceTF
	{\True Số cách chon ra 2 bạn trong 40 bạn lớp 10B là: 780 (cách)}
	{\True Xác suất của biến cố \lq\lq Không bạn nào trong nhóm siêu quậy được gọi\rq\rq\,bằng: $\dfrac{21}{26}$}
	{Xác suất của biến cố \lq\lq Một bạn trong nhóm siêu quậy được gọi\rq\rq\,bằng: $\dfrac{12}{67}$}
	{Xác suất của biến cố \lq\lq Cả hai bạn đươc gọi đều trong nhóm siêu quậy\rq\rq\,bằng: $\dfrac{7}{130}$}
	\loigiai{
		\begin{itemchoice}
			\itemch Đúng. Số cách chọn ra 2 bạn trong 40 bạn lớp 10B là: $\mathrm{C}_{40}^2=780$ (cách).
			\itemch Đúng. Số cách chọn ra 2 bạn trong lớp 10B mà không bạn nào thuộc nhóm siêu quậy là: $\mathrm{C}_{36}^2=630$ (cách). Suy ra $n(A)=630$.
			\itemch Sai. Số cách chọn một bạn trong nhóm siêu quậy là 4 cách. Số cách chọn một bạn
			\itemch Sai. Số cách để cả hai bạn được gọi đều trong nhóm siêu quậy là: $\mathrm{C}_4^2=6$ (cách).
		\end{itemchoice}
	}
\end{ex}
%ex_3
\begin{ex} %[0D0V2-2]
	Hai bạn Nam và Việt, mỗi người gieo môt viên xúc xắc 6 măt cân đối. Khi đó:
	\choiceTF
	{Xác suất để: Nam gieo được số chấm nhỏ hơn $3$ bằng $\dfrac{1}{9}$}
	{\True Xác suất để: Việt gieo được số chấm nhỏ hơn $3$ bằng $\dfrac{1}{3}$}
	{Xác suất để: cả hai bạn đều gieo được số chấm nhỏ hơn $3$ bằng $\dfrac{1}{3}$}
	{\True Xác suất để: cả hai bạn đều gieo được số chấm không nhỏ hơn $4$ bằng $\dfrac{1}{4}$}
	\loigiai{
		\begin{itemchoice}
			\itemch Sai. Không gian mẫu là: $\Omega=\{1 ; 2 ; 3 ; 4 ; 5 ; 6\}$. Do đó, ta có $n(\Omega)=6$.
			\itemch Đúng. Tương tự câu a), ta tính được xác suất để Việt được số chấm nhỏ hơn 3 là $\dfrac{1}{3}$.
			\itemch Sai. Gọi $C$ là biến cố cả hai bạn đều gieo được số chấm nhỏ hơn 3.
			\itemch Đúng. Gọi $D$ là biến cố cả hai bạn đều gieo được số chấm không nhỏ hơn 4.
		\end{itemchoice}
	}
\end{ex}
%ex_4
\begin{ex} %[0D0V2-2]
	Gieo môt con súc sắc. Khi đó:
	\choiceTF
	{\True $n(\Omega)=6$}
	{\True Xác suất để thu đươc măt có số chấm chia hết cho 2 là $\dfrac{1}{2}$}
	{\True Xác suất để thu được măt có số chấm nhỏ hơn 4 là $\dfrac{1}{2}$}
	{Xác suất để thu được măt có số chấm lớn hơn 4 là $\dfrac{1}{2}$}
	loigiai{
		\begin{itemchoice}
			\itemch Đúng. Ta có $\Omega=\{1 ; 2 ; 3 ; 5 ; 6\} \Rightarrow n(\Omega)=6$.
			\itemch Đúng. Gọi $A$ là biến cố: \lq\lq Số chấm thu được chia hết cho 2 \lq\lq .
			\itemch Đúng. Gọi $B$ là biến cố: \lq\lq Số chấm thu được nhỏ hơn 4 \lq\lq .
			\itemch Sai. Gọi $\mathrm{C}$ là biến cố: \lq\lq Số chấm thu được lớn hơn 4\rq\rq.
		\end{itemchoice}
	}
\end{ex}
\Closesolutionfile{ans}
\indapan{3}{ans/ans-0-B1-DS}
\Opensolutionfile{ans}[ans/ans-0-B1-KQ]
\caukq %6cau
\begin{ex}%[0D0H2-7]
	Chọn ngẫu nhiên 2 số trong tập hợp $X=\{1 ; 2 ; 3 ; \ldots ; 50\}$. Tính xác suất của biến cố sau:
	A: \lq\lq Hai số được chọn là số chẵn\rq\rq;\\
	\shortans{$0{,}24$}
	\loigiai{
		Số cách chọn 2 số từ tập hợp $X$ gồm 50 số là: $\mathrm{C}_{50}^2=1225$ (cách).\\
		Do đó, $n(\Omega)=1225$.\\
		Trong tập hợp X có 25 số chẵn $\{2 ; 4 ; 6 \ldots ; 50\}$, nên số cách lấy ra 2 số chẵn là: $\mathrm{C}_{25}^2=300$ (cách). Do đó, $n(A)=300$.\\
		Xác suất của biến cố $A$ là: $\mathrm{P}(A)=\dfrac{n(A)}{n(\Omega)}=\dfrac{300}{1225}=\dfrac{12}{49}\approx0{,}24$.\\
	}
\end{ex}
\begin{ex}%[0D0V2-7]
	Chọn ngẫu nhiên 2 số trong tập hợp $X=\{1 ; 2 ; 3 ; \ldots ; 50\}$. Tính xác suất của biến cố sau: $B$: \lq\lq Trong hai số được chọn có một số lớn hơn 25 , số còn lại nhỏ hơn hoặc bằng 25\rq\rq.
	\shortans{$0{,}51$}
	\loigiai{
		Số cách chọn $2$ số từ tập hợp $X$ gồm $50$ số là: $\mathrm{C}_{50}^2=1225$ (cách).\\
		Số cách chọn một số lớn hơn $25$ là $25$ cách.\\
		Số cách chọn số còn lại nhỏ hơn hoặc bằng $25$ là $25$ cách.\\
		Do đó, ta có $n(B)=25 \cdot 25=625$.\\
		Xác suất của biến cố $B$ là: $\mathrm{P}(A)=\dfrac{n(A)}{n(\Omega)}=\dfrac{625}{1225}=\dfrac{25}{49}\approx0{,}51$.\\
	}
\end{ex}
\begin{ex}%[0D0V2-2]
	Gieo một viên xúc xắc 6 mặt cân đối và đồng chất liên tiếp năm lần. Tính xác suất để mặt 6 chấm xuất hiện ít nhất một lần.
	\shortans{$0{,}6$}
	\loigiai{
		Gọi $A$ là biến cố \lq\lq Mặt 6 chấm không xuất hiện lần nào\rq\rq. Suy ra $A$ là biến cố \lq\lq Mặt 6 chấm xuất hiện ít nhất một lần\rq\rq.\\
		Ta có: $n(\Omega)=6 \cdot 6 \cdot 6 \cdot 6 \cdot 6=7776$, $n(A)=5 \cdot 5 \cdot 5 \cdot 5 \cdot 5=3125$.\\
		Do đó, xác suất của biến cố $A$ là: $\mathrm{P}(A)=\dfrac{n(A)}{n(\Omega)}=\dfrac{3125}{7776}$\\
		Vậy xác suất của biến cố \lq\lq Mặt 6 chấm xuất hiện ít nhất một lần\rq\rq\,là: $\mathrm{P}(A)=1-\dfrac{3125}{7776}=\dfrac{4651}{7776}\approx0{,}6$.
	}
\end{ex}
\begin{ex}%[0D0V2-5]
	Trong tủ có 4 đôi giày khác loại. Bạn Lan lấy ra ngẫu nhiên 2 chiếc giày. Tính xác suất để lấy ra được một đôi giày hoàn chỉnh.
	\shortans{$0{,}14$}
	\loigiai{
		Gọi $A$ là biến cố \lq\lq Lấy ra được một đôi giày hoàn chỉnh\rq\rq.\\
		Ta có: $n(\Omega)=\mathrm{C}_8^2=28, n(A)=4$.\\
		Vậy xác suất của biến cố $A$ là: $\mathrm{P}(A)=\dfrac{n(A)}{n(\Omega)}=\dfrac{4}{28}=\dfrac{1}{7}\approx0{,}14$.
	}
\end{ex}
\begin{ex}%[0D0V2-5]
	Một lớp học có 26 bạn nam và 20 bạn nữ. Chọn ngẫu nhiên một bạn trong lớp. Tính xác suất để bạn được chọn là nam.
	\shortans{$0{,}57$}
	\loigiai{
		Ta có $n(\Omega)=26+20=46$.\\
		Gọi $A$ là biến cố bạn được chọn là nam. Vì lớp học có 26 bạn nam nên có 26 cách chọn một bạn nam. Do đó, ta có $n(A)=26$.\\
		Vậy xác suất của biến cố $A$ là: $\mathrm{P}(A)=\dfrac{n(A)}{n(\Omega)}=\dfrac{26}{46}=\dfrac{13}{23}\approx0{,}57$.
	}
\end{ex}
\begin{ex}%[0D0V2-2]
	Gieo đồng thời hai viên xúc xắc cân đối và đồng chất. Tính xác suất để tổng số chấm xuất hiện trên hai viên xúc xắc bằng 9.
	\shortans{$0{,}11$}
	\loigiai{
		Ta có $n(\Omega)=36$.\\
		Gọi $A$ là biến cố tổng số chấm trên hai viên xúc xắc bằng 9.\\
		$A=\{(3 ; 6),(4 ; 5) ;(5 ; 4) ;(6 ; 3)\}$. Do đó, ta có $n(A)=4$.\\
		Vậy xác suất của biến cố $A$ là: $\mathrm{P}(A)=\dfrac{n(A)}{n(\Omega)}=\dfrac{4}{36}=\dfrac{1}{9}\approx0{,}11$.
	}
\end{ex}
\Closesolutionfile{ans}
\indapan{6}{ans/ans-0-B1-KQ}
\newpage
\subsection{Đề 2}
\Opensolutionfile{ans}[ans/ans-0-B1]
\caulc
\begin{ex}%[0D0N2-5]
	Một chiếc hộp có 9 thẻ được đánh số từ 1 đến 9. Rút ngẫu nhiên hai thẻ rồi nhân hai số ghi trên hai thẻ với nhau. Tính xác suất để kết quả nhận được là số chẵn.
	\choice
	{\True $\dfrac{13}{18}$}
	{$\dfrac{8}{9}$}
	{$\dfrac{4}{9}$}
	{$\dfrac{5}{54}$}
	\loigiai{Rút ngẫu nhiên 2 thẻ trong 9 thẻ, số cách rút là: $\mathrm{C}_9^2=36$.\\
		Để tích nhận được là số lẻ thì phải rút được 2 thẻ đánh số lẻ trong số 5 thẻ đánh số lẻ. Số cách rút là: $\mathrm{C}_5^2=10$\\
		Vậy để tích nhận được là số chẵn thì số cách rút là: $36-10=26$.\\
		Xác suất cần tính là: $P=\dfrac{26}{36}=\dfrac{13}{18}$.
	}
\end{ex}
\begin{ex}%[0D0N2-1]
	Cho $A, B$ là hai biến cố xung khắc. Biết $\mathrm{P}(A)=\dfrac{1}{3}, \mathrm{P}(B)=\dfrac{1}{4}$. Tính $\mathrm{P}(A \cup B)$.
	\choice
	{\True $\dfrac{7}{12}$}
	{$\dfrac{1}{12}$}
	{$\dfrac{1}{7}$}
	{$\dfrac{1}{2}$}
	\loigiai{$$\mathrm{P}(A \cup B)=\mathrm{P}(A)+\mathrm{P}(B)=\dfrac{7}{12}$$
	}
\end{ex}
\begin{ex}%[0D0N2-1]
	Cho $A, B$ là hai biến cố xung khắc. Đẳng thức nào sau đây đúng?
	\choice
	{\True $\mathrm{P}(A \cup B)=\mathrm{P}(A)+\mathrm{P}(B)$}
	{$\mathrm{P}(A \cup B)=\mathrm{P}(A) \cdot \mathrm{P}(B)$}
	{$\mathrm{P}(A \cup B)=\mathrm{P}(A)-\mathrm{P}(B)$}
	{$\mathrm{P}(A \cap B)=\mathrm{P}(A)+\mathrm{P}(B)$}
	\loigiai{Ta có $\mathrm{P}(A \cup B)=\mathrm{P}(A)+\mathrm{P}(B)-\mathrm{P}(A \cap B)$.\\
		Vì $A, B$ là hai biến cố xung khắc nên $A \cap B=\varnothing$. Từ đó suy ra $\mathrm{P}(A \cup B)=\mathrm{P}(A)+\mathrm{P}(B)$.
	}
\end{ex}
\begin{ex}%[0D0N2-1]
	Gọi $A$ và $B$ là hai biến cố liên quan đến phép thử ngẫu nhiên $T$. Cho $\mathrm{P}(A)=\dfrac{1}{4}, \mathrm{P}(A \cup B)=\dfrac{1}{2}$. Biết $A, B$ là hai biến cố xung khắc, thì $\mathrm{P}(B)$ bằng:
	\choice
	{$\dfrac{3}{4}$}
	{$\dfrac{1}{8}$}
	{$\dfrac{1}{3}$}
	{\True $\dfrac{1}{4}$}
	\loigiai{Ta có: $A, B$ là hai biến cố xung khắc nên\\
		$$\mathrm{P}(A \cup B)=\mathrm{P}(A)+\mathrm{P}(B) \Rightarrow \mathrm{P}(B)=\mathrm{P}(A \cup B)-\mathrm{P}(A)=\dfrac{1}{4}.$$
	}
\end{ex}
\begin{ex}%[0D0N2-1]
	Cho hai biến cố xung khắc $A$ và $B$. Biết $\mathrm{P}(A)=\dfrac{1}{4}, \mathrm{P}(A \cup B)=\dfrac{1}{2}$. Tính $\mathrm{P}(B)$.
	\choice
	{$\dfrac{1}{3}$}
	{$\dfrac{1}{8}$}
	{\True $\dfrac{1}{4}$}
	{$\dfrac{3}{4}$}
	\loigiai{
		Vì $A$ và $B$ là hai biến cố xung khắc nên $A \cap B= \varnothing $.\\
		Khi đó ta có: $\mathrm{P}(A \cup B)=\dfrac{1}{2} \Leftrightarrow \mathrm{P}(A)+\mathrm{P}(B)=\dfrac{1}{2} \Leftrightarrow \mathrm{P}(B)=\dfrac{1}{2}-\dfrac{1}{4}=\dfrac{1}{4}$.
	}
\end{ex}
\begin{ex} %[0D0N2-1]
	Cho hai biến cố $A$ và $B$ có $\mathrm{P}(A)=\dfrac{1}{3}$, $\mathrm{P}(B)=\dfrac{1}{4}$, $\mathrm{P}(A \cup B)=\dfrac{1}{2}$. Ta kết luận hai biến cố  $A$ và $B$ là 
	\choice
	{Độc lập}
	{Không độc lập}
	{Xung khắc}
	{\True Không xung khắc}
	\loigiai{
		Với hai biến cố $A$ và $B$ bất kì liên quan đến cùng một phép thử, ta có: \\
		$\mathrm{P}(A \cup B)=\mathrm{P}(A)+\mathrm{P}(B)-\mathrm{P}(A \cap B)$  \\
		$\Rightarrow \mathrm{P}(A \cap B)=\mathrm{P}(A)+\mathrm{P}(B)-\mathrm{P}(A \cup B)=\dfrac{1}{3}+\dfrac{1}{4}-\dfrac{1}{2}=\dfrac{1}{12} \neq 0$\\ 
		$\Rightarrow A \cap B \neq \varnothing \Rightarrow A$ và $B$ không xung khắc.
	}
\end{ex}
\begin{ex}%[0D0N2-1]
	Gieo một đồng tiền liên tiếp 3 lần. Gọi $A$ là biến cố ít nhất một lần xuất hiện mặt sấp. Tính xác suất $\mathrm{P}(A)$ của biến cố $A$.
	\choice
	{$\mathrm{P}(A)=\dfrac{3}{8}$}
	{$\mathrm{P}(A)=\dfrac{1}{4}$}
	{$\mathrm{P}(A)=\dfrac{1}{2}$}
	{\True $\mathrm{P}(A)=\dfrac{7}{8}$}
	\loigiai{Không gian mẫu là: $\Omega=\{S S S, S N N, N S N, N N S, S S N, S N S, N S S, N N N\}$.\\
		$\Rightarrow n(\Omega)=8$.\\
		$A$ là biến cố ít nhất một lần xuất hiện mặt sấp nên $\overline{A}$ là biến cố không lần nào xuất hiện mặt sấp. Ta có $\overline{A}=\{N N N\} \Rightarrow n(\overline{A})=1$.\\
		Xác suất của biến cố $\overline{A}$ là: $\mathrm{P}(\overline{A})=\dfrac{n(\overline{A})}{n(\Omega)}=\dfrac{1}{8}$.\\
		Xác suất của biến cố $A$ là: $\mathrm{P}(A)=1-\mathrm{P}(\overline{A})=1-\dfrac{1}{8}=\dfrac{7}{8}$.
	}
\end{ex}
\begin{ex}%[0D0N2-1]
	Do con xúc xắc không cân đối nên mỗi khi gieo con xúc xắc đó thì xác suất xuất hiện mặt 2 chấm và mặt 5 chấm lần lượt là $0{,}35$ và $0{,}25$ . Xác suất xuất hiện mặt 2 chấm hoặc mặt 5 chấm là
	\choice
	{\True $0{,}6$}
	{$0{,}0875$}
	{$0{,}4875$}
	{$0{,}9125$}
	\loigiai{
		Phép thử T: \lq\lq Gieo một con xúc xắc không cân đối\rq\rq.\\
		Gọi $A$ là biến cố: \lq\lq xuất hiện mặt 2 chấm\rq\rq;\\
		Gọi $B$ là biến cố: \lq\lq xuất hiện mặt 5 chấm\rq\rq;\\
		Gọi $C$ là biến cố: \lq\lq xuất hiện mặt 2 chấm hoặc mặt 5 chấm\rq\rq.\\
		Từ giả thiết ta có $\mathrm{P}(A)=0,35$, $\mathrm{P}(B)=0{,}25$.\\
		Mà $C= A\cup B$ và $A \cap B =\varnothing$ nên xác suất cần tìm là $\mathrm{P}(C)=\mathrm{P}(A)+\mathrm{P}(B)=0{,}6$.
	}
\end{ex}
\begin{ex}%[0D0N2-1]
	Cho $A, B$ là hai biến cố xung khắc. Biết $\mathrm{P}(A)=\dfrac{1}{5}, \mathrm{P}(A \cup B)=\dfrac{1}{3}$. Khi đó $\mathrm{P}(B)$ là
	\choice
	{$\dfrac{3}{5}$}
	{$\dfrac{8}{15}$}
	{\True $\dfrac{2}{15}$}
	{$\dfrac{1}{15}$}
	\loigiai{Vì $A, B$ là hai biến cố xung khắc nên ta có:\\
		$$\mathrm{P}(A \cup B)=\mathrm{P}(A)+\mathrm{P}(B) \Rightarrow \mathrm{P}(B)=\mathrm{P}(A \cup B)-\mathrm{P}(A)=\dfrac{2}{15}$$\\
	}
\end{ex}
\begin{ex}%[0D0N2-1]
	Một nhóm gồm 6 học sinh nam và 4 học sinh nữ. Chọn ngẫu nhiên đồng thời 3 học sinh trong nhóm đó. Xác suất để trong 3 học sinh được chọn luôn có học sinh nữ bằng
	\choice
	{\True $\dfrac{5}{6}$}
	{$\dfrac{2}{3}$}
	{$\dfrac{1}{6}$}
	{$\dfrac{1}{3}$}
	\loigiai{Số phần từ của không gian mẫu $n(\Omega)=\mathrm{C}_{10}^3=120$.\\
		Gọi $A$ là biến cố sao cho 3 học sinh được chọn có học sinh nũ,\\
		$\Rightarrow \overline{A}$ là biến cố sao cho 3 học sinh được chọn không có học sinh nữ $\Rightarrow n(\overline{A})=\mathrm{C}_6^3=20$.\\
		Vậy xác suất cần tìm $\mathrm{P}(A)=1-\mathrm{P}(\overline{A})=1-\dfrac{n(\overline{A})}{n(\Omega)}=\dfrac{5}{6}$.\\
	}
\end{ex}
\begin{ex}%[0D0N2-1]
	Bạn Trang có 10 đôi tất. Trong lúc vội Trang đã lấy ngẫu nhiên 4 chiếc tất. Tính xác suất để trong 4 chiếc tất có ít nhất một đôi tất.
	\choice
	{$\dfrac{224}{323}$}
	{\True $\dfrac{99}{323}$}
	{$\dfrac{204}{323}$}
	{$\dfrac{9}{323}$}
	\loigiai{Số phần tử không gian mẫu: $|\Omega|=\mathrm{C}_{20}^4$.\\
		Gọi $A$ là biến cố: \lq\lq Có ít nhất một đôi tất trong 4 chiếc tất được lấy ngẫu nhiên từ 10 đôi tất\rq\rq.\\
		Suy ra $\overline{A}$ là biến cố: \lq\lq Không có đôi tất nào trong 4 chiếc tất được lấy ngẫu nhiên từ 10 đôi tất\rq\rq.\\
		$|\overline{A}|=\mathrm{C}_{10}^4 \cdot 2^4$ Chọn 4đôi tất và lấy mỗi đôi 1chiếc).\\
		$$\mathrm{P}(A)=1-\mathrm{P}(\overline{A})=1-\dfrac{\mathrm{C}_{10}^4 \cdot 2^4}{\mathrm{C}_{20}^4}=\dfrac{99}{323}$$\\
	}
\end{ex}
\begin{ex}%[0D0N2-1]
	Cho hai biến cố độc lập $A$ và $B$. Biết $\mathrm{P}(A)=\dfrac{1}{4}, \mathrm{P}(A \cup B)=\dfrac{1}{2}$. Tính $\mathrm{P}(B)$.
	\choice
	{$\dfrac{3}{4}$}
	{$\dfrac{1}{8}$}
	{$\dfrac{1}{4}$}
	{\True $\dfrac{1}{3}$}
	\loigiai{Ta có: $\mathrm{P}(A \cup B)=\mathrm{P}(A)+\mathrm{P}(B)-\mathrm{P}(A \cap B)$\\
		Mà $A$ và $B$ là hai biến cố độc lập nên $\mathrm{P}(A \cup B)=\mathrm{P}(A)+\mathrm{P}(B)-\mathrm{P}(A) \cdot \mathrm{P}(B)$\\
		$\Rightarrow \dfrac{1}{2}=\dfrac{1}{4}+\mathrm{P}(B)-\dfrac{1}{4} \mathrm{P}(B) \Rightarrow \mathrm{P}(B)=\dfrac{1}{3}$
	}
\end{ex}
\Closesolutionfile{ans}
\indapan{6}{ans/ans-0-B1}
\cauds
\Opensolutionfile{ans}[ans/ans-0-B1-DS]
\begin{ex} %[0D0V2-5]
	Trong hộp có chứa 7 bi xanh, 5 bi đỏ, 2 bi vàng có kích thước và khối lượng như nhau. Lấy ngẫu nhiên từ trong hộp 6 viên bi. Khi đó:
	\choiceTF
	{\True Xác suất để có đúng một màu bằng: $\dfrac{1}{429}$}
	{\True Xác suất để có đúng hai màu đỏ và vàng bằng: $\dfrac{1}{429}$}
	{Xác suất để có ít nhất 1 bi đỏ bằng: $\dfrac{139}{143}$}
	{Xác suất để có ít nhất 2 bi xanh bằng: $\dfrac{32}{39}$}
	\loigiai{Chọn ngẫu nhiên 6 viên bi trong 14 viên bi, có $\mathrm{C}_{14}^6$ cách.\\
		Vậy số phần tử của không gian mẫu $n(\Omega)=\mathrm{C}_{14}^6=3003$
		\begin{itemchoice}
			\itemch Đúng. Gọi A: \lq\lq 6 viên được chọn có đúng một màu\rq\rq.\\
			$n(A)=\mathrm{C}_7^6$. Suy ra $\mathrm{P}(A)=\dfrac{n(A)}{n(\Omega)}=\dfrac{\mathrm{C}_7^6}{\mathrm{C}_{14}^6}=\dfrac{1}{429}$.
			\itemch Đúng. Gọi biến cố B: \lq\lq 6 viên được chọn có đúng hai màu đỏ và vàng\rq\rq.\\
			Số trường hợp thuận lợi cho $B$ là:\\
			Trường hợp 1: Chọn được 1 vàng và 5 đỏ, có $\mathrm{C}_2^1 \cdot \mathrm{C}_5^5=2$ cách.\\
			Trường hợp 2: Chọn được 2 vàng và 4 đỏ, có $\mathrm{C}_2^2 \cdot \mathrm{C}_5^4=5$ cách.\\
			$n(B)=2+5=7$. Suy ra $\mathrm{P}(B)=\dfrac{n(B)}{n(\Omega)}=\dfrac{7}{\mathrm{C}_{14}^6}=\dfrac{1}{429}$
			\itemch Sai. Gọi $C$: \lq\lq 6 viên được chọn có ít nhất 1 bi đỏ\rq\rq.\\
			Biến cố đối $\overline{C}$ : \lq\lq Tất cả 6 viên được chọn đều không có bi đỏ\rq\rq.\\
			$n(\overline{C})=\mathrm{C}_9^6=84$. Suy ra $\mathrm{P}(\overline{C})=\dfrac{n(\overline{C})}{n(\Omega)}=\dfrac{4}{143}$.
			$$
			\mathrm{P}(C)+\mathrm{P}(\overline{C})=1 \Rightarrow \mathrm{P}(C)=1-\mathrm{P}(\overline{C})=\dfrac{139}{143}
			$$
			\itemch Sai. Gọi biến cố D: \lq\lq 6 viên được chọn có ít nhất 2 bi xanh\rq\rq.\\
			Biến cố đối $\overline{D}$ : \lq\lq 6 viên được chọn có nhiều nhất 1 bi xanh\rq\rq.\\
			Số trường hợp thuận lợi cho $\overline{D}$ là:\\
			Trường hợp 1: Chọn được 6 bi đỏ, vàng, có $\mathrm{C}_7^6=7$ cách.\\
			Trường hợp 2: Chọn được 1 bi xanh và 5 bi đỏ, vàng, có $\mathrm{C}_7^1 \cdot \mathrm{C}_7^5=147$ cách.\\
			$n(\overline{D})=7+147=154$. \\
			Suy ra $\mathrm{P}(\overline{D})=\dfrac{n(\overline{D})}{n(\Omega)}=\dfrac{2}{39}$.
			$\mathrm{P}(D)+\mathrm{P}(\overline{D})=1 \Rightarrow \mathrm{P}(D)=1-\mathrm{P}(\overline{D})=\dfrac{37}{39}$
		\end{itemchoice}
	}
\end{ex}
%ex_11
\begin{ex} %[0D0V2-2]
	Gieo hai con xúc xắc. Khi đó:
	\choiceTF
	{\True Xác suất \lq\lq Số chấm xuất hiện trên hai con xúc xắc hơn kém nhau 2 chấm\rq\rq\,  bằng: $\dfrac{2}{9}$}
	{\True Xác suất \lq\lq Tích số chấm xuất hiện trên hai con xúc xắc chia hết cho 5 \lq\lq  bằng: $\dfrac{11}{36}$}
	{Xác suất \lq\lq Tích số chấm xuất hiện trên hai con xúc xắc là một số chẵn\rq\rq\,  bằng: $\dfrac{5}{6}$}
	{\True Xác suất \lq\lq Tổng số chấm xuất hiện trên hai con xúc xắc là số lẻ\rq\rq\,  bằng: $\dfrac{1}{2}$}
	\loigiai{
		Không gian mẫu $\Omega=\{(i ; j) \mid i, j=1,2,3, \ldots, 6\}$.
		Số các kết quả có thể xảy ra của phép thử là $n(\Omega)=6.6=36$.
		\begin{itemchoice}
			\itemch Đúng. Gọi $A$ là biến cố \lq\lq Số chấm xuất hiện trên hai con xúc xắc hơn kém nhau 2 chấm\rq\rq.\\
			$A=\{(1 ; 3) ;(2 ; 4) ;(3 ; 5) ;(4 ; 6) ;(3 ; 1) ;(4 ; 2) ;(5 ; 3) ;(6 ; 4)\}$, $n(A)=8$. Suy ra\\ $\mathrm{P}(A)=\dfrac{n(A)}{n(\Omega)}=\dfrac{2}{9}$.
			\itemch Đúng. Gọi $B$ là biến cố \lq\lq Tích số chấm xuất hiện trên hai con xúc xắc chia hết cho 5\rq\rq.\\
			$B=\{(1 ; 5) ;(2 ; 5) ;(3 ; 5) ;(4 ; 5) ;(5 ; 5) ;(6 ; 5) ;(5 ; 1) ;(5 ; 2) ;(5 ; 3) ;(5 ; 4) ;(5 ; 6)\} \Rightarrow n(B)=11$.\\
			Suy ra $\mathrm{P}(B)=\dfrac{n(B)}{n(\Omega)}=\dfrac{11}{36}$.
			\itemch Sai. Gọi C: \lq\lq Tích số chấm xuất hiện trên hai con xúc xắc là một số chẵn\rq\rq.\\ Biến cố đối $\overline{C}$ : \lq\lq Tích số chấm xuất hiện trên hai con xúc xắc là một số lẻ\rq\rq.\\ $n(\overline{C})=3.3=9$. Suy ra $\mathrm{P}(\overline{C})=\dfrac{n(\overline{C})}{n(\Omega)}=\dfrac{1}{4}$.
			$\mathrm{P}(C)+\mathrm{P}(\overline{C})=1 \Rightarrow \mathrm{P}(C)=1-\mathrm{P}(\overline{C})=\dfrac{3}{4}$
			\itemch Đúng Gọi D: \lq\lq Tổng số chấm xuất hiện trên hai con xúc xắc là số lẻ\rq\rq.\\
			$\overline{D}$ : \lq\lq Tổng số chấm xuất hiện trên hai con xúc xắc là số chẵn\rq\rq.\\
			Ta có tổng số chấm xuất hiện trên hai con xúc xắc là số chẵn khi và chỉ khi cả hai số đó đều là số lẻ hoặc đều là số chẵn.\\
			$n(\overline{D})=2.3 .3=18 \text {. Suy ra } \mathrm{P}(\overline{D})=\dfrac{n(\overline{D})}{n(\Omega)}=\dfrac{1}{2} \Rightarrow P(D)=1-\mathrm{P}(\overline{D})=\dfrac{1}{2}$
		\end{itemchoice}
	}
\end{ex}
\begin{ex} %[0D0V2-5]
	Một hộp đựng 9 thẻ được đánh số từ 1 đến 9 . Rút ngẫu nhiên 5 thẻ. Khi đó:
	\choiceTF
	{\True Xác suất \lq\lq Các thẻ ghi số $1,2,3$ được rút\rq\rq\,  bằng: $\dfrac{5}{42}$}
	{Xác suất \lq\lq Có đúng 1 trong 3 thẻ ghi số 1,2,3 được rút\rq\rq\,  bằng: $\dfrac{6}{11}$}
	{\True Xác suất \lq\lq Không thẻ nào trong 3 thẻ ghi số 1,2,3 được rút\rq\rq\,  bằng: $\dfrac{1}{21}$}
	{\True Xác suất \lq\lq Có ít nhất một trong 3 thẻ ghi số 1,2,3 được rút\rq\rq\,  bằng: $\dfrac{20}{21}$}
	\loigiai{
		$n(\Omega)=\mathrm{C}_9^5=126$.
		\begin{itemchoice}
			\itemch Đúng. $n(A)=\mathrm{C}_6^2=15 ; \mathrm{P}(A)=\dfrac{n(A)}{n(\Omega)}=\dfrac{15}{126}=\dfrac{5}{42}$.
			\itemch Sai. $n(B)=\mathrm{C}_3^1 \mathrm{C}_6^4=45 ; \mathrm{P}(B)=\dfrac{n(B)}{n(\Omega)}=\dfrac{45}{126}=\dfrac{5}{14}$.
			\itemch Đúng. $n(C)=\mathrm{C}_6^5=6 ; \mathrm{P}(C)=\dfrac{n(C)}{n(\Omega)}=\dfrac{6}{126}=\dfrac{1}{21}$.
			\itemch Đúng. Biến cố đối $\overline{D}=C$ : \lq\lq Không thẻ nào trong 3 thẻ ghi số $1,2,3$ được rút\rq\rq.\\
			$\Rightarrow \mathrm{P}(\overline{D})=\dfrac{1}{21} ; \mathrm{P}(D)=1-\mathrm{P}(\overline{D})=\dfrac{20}{21}$
		\end{itemchoice}
	}
\end{ex}
\begin{ex} %[0D0V2-5]
	Một bình chứa 16 viên bi với 7 viên bi trắng, 6 viên bi đen và 3 viên bi đỏ. Lấy ngẫu nhiên 3 viên bi. Khi đó:
	\choiceTF
	{\True Số phần tử của không gian mẫu là $560$.}
	{\True Xác suất lấy được cả 3 viên bi đỏ, bằng: $\dfrac{1}{560}$.}
	{Xác suất lấy được cả 3 viên bi không đỏ, bằng: $\dfrac{43}{280}$}
	{\True Xác suất lấy được cả 1 viên bi trắng, 1 viên bi đen, 1 viên bi đỏ, bằng: $\dfrac{9}{40}$}
	\loigiai{
		\begin{itemchoice}
			\itemch Đúng. Không gian mẫu $n(\Omega)=\mathrm{C}_{16}^3=560$.
			\itemch Đúng. Gọi $A$ là biến cố: \lq\lq lấy được 3 viên bi đỏ\rq\rq.\\
			Ta có $n(A)=1$. Vậy $\mathrm{P}(A)=\dfrac{1}{560}$.
			\itemch Sai. Gọi $A$ : \lq\lq lấy được 3 viên bi không đỏ\rq\rq\,  thì $A$ : \lq\lq lấy được 3 viên bi trắng hoặc đen\rq\rq. Có $7+6=13$ viên bi trắng hoặc đen. Ta có $n(A)=\mathrm{C}_{13}^3=286$.\\
			Vậy $\mathrm{P}(A)=\dfrac{286}{560}=\dfrac{143}{280}$.
			\itemch Đúng. Gọi $A$ : \lq\lq lấy được 1 viên bi trắng, 1 viên vi đen, 1 viên bi đỏ\rq\rq.\\
			Ta có: $n(A)=7.6 \cdot 3=126$. Vậy $\mathrm{P}(A)=\dfrac{126}{560}=\dfrac{9}{40}$.
		\end{itemchoice}
	}
\end{ex}
\Closesolutionfile{ans}
\indapan{3}{ans/ans-0-B1-DS}
\Opensolutionfile{ans}[ans/ans-0-B1-KQ]
\caukq
\begin{ex}
	Từ bộ bài tây gồm 52 quân bài, người ta rút ra ngẫu nhiên 2 quân bài. Tính xác suất để rút được 2 quân bài khác màu.
	\shortans{$0{,}51$}
	\loigiai{
		Số cách để rút ra ngẫu nhiên 2 quân bài từ bộ bài tây gồm 52 quân bài mà không quan trọng thứ tự là: $\mathrm{C}_{52}^2=1326$ (cách). Do đó, ta có $n(\Omega)=1326$.\\
		Gọi $A$ là biến cố rút được hai quân bài khác màu.\\
		Vì bộ bài tây gồm 26 quân bài đỏ và 26 quân bài đen nên số cách rút được hai quân bài khác màu là: $\mathrm{C}_{26}^1 \cdot \mathrm{C}_{26}^1=676$ (cách). Do đó, ta có $n(A)=676$.\\
		Vậy xác suất của biến cố $A$ là: $\mathrm{P}(A)=\dfrac{n(A)}{n(\Omega)}=\dfrac{676}{1326}=\dfrac{26}{51}\approx0{,}51$.\\
	}
\end{ex}
\begin{ex} 
	Một lô hàng có 14 sản phẩm, trong đó có đúng 2 phế phẩm. Chọn ngẫu nhiên 8 sẩn phẩm trong lô hàng đó. Tính xác suất của biến cố \lq\lq Trong 8 sản phẩm được chọn có không quá 1 phế phẩm\rq\rq.
	\shortans{$0{,}69$}
	\loigiai{
		Số cách chọn ngẫu nhiên 8 sản phẩm là: $\mathrm{C}_{14}^8=3003$.\\
		Số cách chọn ngẫu nhiên 8 sản phẩm mà không có phế phẩm là: $\mathrm{C}_{12}^8=495$.\\
		Số cách chọn ngẫu nhiên 8 sản phẩm mà trong đó có đúng 1 phế phẩm là: $\mathrm{C}_2^1 \cdot \mathrm{C}_{12}^7=1584$.\\
		Suy ra số cách chọn ngẫu nhiên 8 sản phẩm mà trong đó có không quá 1 phế phẩm là: $495+1584=2079$.\\
		Vậy xác suất của biến cố \lq\lq Trong 8 sản phẩm được chọn có không quá 1 phế phẩm\rq\rq\,  là: $\dfrac{2079}{3003}=\dfrac{9}{13}\approx0{,}69$.\\
	}
\end{ex}
\begin{ex}
	Xếp ngẫu nhiên 4 bạn nam và 4 bạn nữ thành một hàng dọc. Tính xác suất của biến cố \lq\lq Xếp được các bạn nam và bạn nữ đứng xen kẽ nhau\rq\rq.
	\shortans{$0{,}03$}
	\loigiai{
		Giả sử các vị trí của hàng dọc được đánh số thứ tự từ đầu hàng là $1,2,3,4,5,6,7,8$. Số cách xếp 8 bạn thành một hàng dọc là $8!=40320$.\\
		Xếp các bạn nam và bạn nữ đứng xen kẽ nhau có hai trường hợp:\\
		Truờng hợp 1: Các bạn nam đứng ở các vị trí số lẻ còn các bạn nữ đứng ở các vị trí số chẵn. Số cách xếp như vậy là $4!.4!=576$.\\
		Truờng hợp 2: Các bạn nữ đứng ở các vị trí số lẻ còn các bạn nam đứng ở các vị trí số chẵn. Số cách xếp như vậy là $4!\cdot 4!=576$.\\
		Vậy xác suất của biến cố \lq\lq Xếp được các bạn nam và bạn nữ đứng xen kẽ nhau\rq\rq\,  là: $\dfrac{576+576}{40320}=\dfrac{1}{35}\approx0{,}03$.\\
	}
\end{ex}
\begin{ex}Cho một đa giác đều 12 đỉnh nội tiếp đường tròn. Chọn ngẫu nhiên 4 đỉnh của đa giác, tính xác suất để 4 đỉnh được chọn ra tạo thành một hình chữ nhật.
	\shortans{$0{,}03$}
	\loigiai{
		Số phần tử không gian mẫu là $n(\Omega)=\mathrm{C}_{12}^4$.\\
		Gọi $A$ là biến cố \lq\lq Chọn được 4 đỉnh tạo thành hình chữ nhật\rq\rq.\\
		Đa giác đều đã cho có $\dfrac{12}{2}=6$ đường chéo lớn.\\
		Mỗi hình chữ nhật được chọn phải có 2 trong 6 đường chéo trên.\\
		Vì vậy $n(A)=\mathrm{C}_6^2$. Suy ra $\mathrm{P}(A)=\dfrac{n(A)}{n(\Omega)}=\dfrac{\mathrm{C}_6^2}{\mathrm{C}_{12}^4}=\dfrac{1}{33}\approx0{,}03$.
	}
\end{ex}
\begin{ex}
	Cho đa giác có 12 đỉnh. Chọn ngẫu nhiên 3 đỉnh của đa giác đó. Tìm xác suất để 3 đỉnh được chọn tạo thành một tam giác không có cạnh nào là cạnh của đa giác đã cho.
	\shortans{$0{,}51$}
	\loigiai{
		Số tam giác được tạo từ 3 đỉnh trong 12 đỉnh là $\mathrm{C}_{12}^3$. Vì vậy $n(\Omega)=\mathrm{C}_{12}^3$.\\
		Gọi $A$ : \lq\lq Chọn được tam giác không có cạnh nào là cạnh của đa giác\rq\rq.\\
		Xét số tam giác có 2 cạnh là cạnh của đa giác:\\
		Các tam giác này sẽ có 3 đỉnh là 3 đỉnh liên tiếp của đa giác tức là có 2 cạnh là 2 cạnh liên tiếp của đa giác, 2 cạnh này cắt nhau tại 1 đỉnh, mà đa giác này có 12 đỉnh nên có 12 tam giác thỏa mãn trường hợp này.\\
		Xét số tam giác có 1 cạnh là cạnh của đa giác:\\
		Trước tiên ta chọn 1 cạnh trong 12 cạnh của đa giác nên có 12 cách chọn.\\
		Tiếp theo chọn 1 đỉnh còn lại trong 8 đỉnh (trừ 2 đỉnh tạo nên cạnh đã chọn và 2 đỉnh liền kể với cạnh đã chọn).\\
		Do đó trong trường hợp này có 8.12 tam giác.\\
		Số tam giác có ít nhất một cạnh là cạnh của đa giác là $12+8.12=108$.\\
		Suy ra: $n(A)=\mathrm{C}_{12}^3-108=112$. Vậy $\mathrm{P}(A)=\dfrac{n(A)}{n(\Omega)}=\dfrac{28}{55}\approx0{,}51$.
	}
\end{ex}
\begin{ex}Từ 12 học sinh gồm 5 học sinh giỏi, 4 học sinh khá, 3 học sinh trung bình, giáo viên muốn thành lập 4 nhóm làm 4 bài tập lớn khác nhau, mỗi nhóm 3 học sinh. Tính xác suất để nhóm nào cũng có học sinh giỏi và học sinh khá.
	\\
	\shortans{$0{,}09$}
	\loigiai{
		\\
		Ta có số phần tử không gian mẫu là $n(\Omega)=\mathrm{C}_{12}^3 \cdot \mathrm{C}_9^3 \cdot \mathrm{C}_6^3 \cdot \mathrm{C}_3^3$.\\
		Gọi $A$ là biến cố: \lq\lq Để nhóm nào cũng có học sinh giỏi và học sinh khá\rq\rq.\\
		Đánh số 4 nhóm là $A, B, C, D$.\\
		Bước 1: Xếp vào mỗi nhóm một học sinh khá có 4! cách.\\
		Bước 2: Xếp 5 học sinh giỏi vào 4 nhóm thì có 1 nhóm có 2 học sinh giỏi. Chọn nhóm có 2 học sinh giỏi có 4 cách, chọn 2 học sinh giỏi có $\mathrm{C}_5^2$ cách, xếp 3 học sinh giỏi còn lại có 3 ! cách.\\
		Bước 3: Xếp 3 học sinh trung bình có 3! cách.
		$$
		\mathrm{P}(A)=\dfrac{4!\cdot 4 \cdot \mathrm{C}_5^2 \cdot 3!\cdot 3!}{\mathrm{C}_{12}^3 \mathrm{C}_9^3 \mathrm{C}_6^3 \mathrm{C}_3^3}=\dfrac{36}{385}\approx0{,09}$$.
	}
\end{ex}
\Closesolutionfile{ans}
\indapan{6}{ans/ans-0-B1-KQ}
